\section{Experiments}
\label{sec:experiments}

\subsection{Setup}
\label{sec:setup}
\paragraph{Data.} The stochastic core is German Credit (UCI, $d{=}\nGerDim$ attributes, 1{,}000 rows, 700/300 stratified split) and ACSEmployment from folktables \citep{ding2021retiring} ($d{=}\nAcsDim$; California, 2018 one-year person file, 40{,}000 rows subsampled with a fixed seed, 28{,}000/12{,}000 split). The deterministic-regime controls are ACSIncome ($d{=}10$) and Pima diabetes ($d{=}8$). Sources, versions and retrieval dates are in the reproduction notes.

\paragraph{Models.} Logistic regression (LR), XGBoost (depth 4, 200 trees), a random forest (RF; 100 trees, depth $\le8$) and an MLP (two hidden layers of 64 units), trained once with a fixed seed. Held-out AUC is $\nGerAucMin$--$\nGerAucMax$ on German Credit and $\nAcsAucMin$--$\nAcsAucMax$ on ACSEmployment (appendix table). Categorical attributes are ordinal-coded for trees; LR and MLP one-hot encode inside the pipeline so that masking an attribute masks its whole group. The explained value function is the log-odds for LR, XGBoost and MLP and the probability for RF, for which interventional TreeSHAP is exact; results are therefore not pooled across the RF and the other models. A sensitivity check with the other value function is in Appendix~\ref{app:valfn}.

\paragraph{Explainers, budget and targets.} Kernel SHAP (\texttt{shap} 0.51.0 \texttt{KernelExplainer}) is the main explainer; \texttt{PermutationExplainer} is the second. The background is 50 training rows, fixed per dataset. The budget is counted in model evaluations, $B=\texttt{nsamples}\times50$; the library default is $\texttt{nsamples}=2d+2048$ (e.g.\ $B{=}\nGerOneEvals$ on German Credit, about $\nGerOneSec$\,s per explanation on the shared CPU host used). We study $\nicefrac14$, $1$ and $4$ times the default for all cells and $16\times$ on German Credit for 20 instances; the $16\times$ ACSEmployment cell was cut for throughput. $\ell_1$ selection is off in the primary arm; the library default is run separately at $1\times$. Exact targets use the same background and value function: closed form for LR, interventional TreeSHAP for XGBoost and RF, full enumeration for all ACS models, and enumeration on the first 20 instances for the German MLP (so its accuracy rows use 20 instances). Targets were cross-checked against enumeration (agreement to $2\times10^{-16}$ for LR, $10^{-6}$ for XGBoost and $2\times10^{-4}$ for RF).

\paragraph{Instances, runs and statistics.} Fifty instances per dataset, shared by all models: 25 random test rows and 25 test rows nearest the decision boundary among a random pool of 300. Each cell has $R{=}20$ independent explainer seeds. Pairwise metrics use all 190 run pairs per instance. Intervals are 95\% nested-bootstrap percentile intervals (1{,}000 draws: instances, then runs within instance); the instance is the unit of inference. Total compute was about 26 CPU-hours on a host oversubscribed during the runs. The null reference for agreement is a random permutation of the features of the same vector (chance level).

\subsection{Single runs at the library default budget}
\label{sec:main}
\begin{figure}[t]
\centering
\includegraphics[width=\linewidth]{fig1_budget_curves.pdf}
\caption{Single-run reliability of Kernel SHAP against budget on the two stochastic datasets (columns). The budget is the multiple of the library default $\texttt{nsamples}=2d+2048$; $B=\texttt{nsamples}\times50$ model evaluations is printed under each tick. Top row: N1, mean top-3 Jaccard between two runs. Middle: N2, fraction of explanations whose top-3 set equals the exact target's. Bottom: N2, median relative $\ell_2$ error (log scale; triangles: below $10^{-4}$, i.e.\ exact to numerical noise, as for the additive LR). Points are means over 50 instances (German MLP error rows: 20 instances) of per-instance means over runs; bars are 95\% nested-bootstrap intervals. Dashed hollow: permutation SHAP at $1\times$ and $4\times$; crosses: the \texttt{shap} default $\ell_1$ setting at $1\times$; dotted: permuted-feature chance level; dash-dotted: thresholds of the decision rule. $16\times$ exists for German Credit only.}
\label{fig:budget}
\end{figure}

\paragraph{With $\ell_1$ off, one run is repeatable and accurate for top-feature identification.} At $1\times$ the mean top-3 hit rate against the exact target (per model, over 50 instances) is $\nGerOneTopHitEstMin$--$\nGerOneTopHitEstMax$ on German Credit (lowest 95\% lower bound $\nGerOneTopHitLoMin$) and $\nAcsOneTopHitEstMin$--$\nAcsOneTopHitEstMax$ on ACSEmployment; the median relative $\ell_2$ error is at most $\nGerOneRelErrEstMax$ and $\nAcsOneRelErrEstMax$ respectively, and the mean pairwise Spearman correlation of $|\phi|$ rankings between runs is $\ge\nGerOneSpearEstMin$ and $\ge\nAcsOneSpearEstMin$. LR is exact at machine precision at every budget (the model is additive), so its rows are a check on the metrics, not evidence about sampling noise; XGBoost and RF have an exact algorithm that one would use instead of sampling. The MLP is therefore the only model here for which sampling is the practical option: its top-3 hit rate at $1\times$ is $\nGerMlpOneTopHitEst$ [$\nGerMlpOneTopHitLo$, $\nGerMlpOneTopHitHi$] on German Credit (20 instances, exact targets by enumeration) and $\nAcsMlpOneTopHitEst$ [$\nAcsMlpOneTopHitLo$, $\nAcsMlpOneTopHitHi$] on ACSEmployment, with median relative error $\nGerMlpOneRelErrEst$ and $\nAcsMlpOneRelErrEst$. All $\nGerOneVerRel$ German Credit and $\nAcsOneVerRel$ ACSEmployment cells at $1\times$ meet the reliability rule (Table~\ref{tab:verdicts}); so do the $\nicefrac14\times$ cells (top-3 hit $\ge\nGerQTopHitEstMin$, relative error $\le\nGerQRelErrEstMax$ on German Credit), all $4\times$ and $16\times$ cells, and permutation SHAP at $1\times$ and $4\times$ (Table~\ref{tab:verdicts}, Figure~\ref{fig:budget}). Error falls with budget as expected (Figure~\ref{fig:budget}). On the controls, the library default enumerates all coalitions, so repeatability is exact by construction and these cells are not evidence of reliability for larger $d$.

\paragraph{The shap default $\ell_1$ setting is the exception.} In \texttt{shap} 0.51.0 the default of \texttt{KernelExplainer.shap\_values} is \texttt{l1\_reg='num\_features(10)'}, and the selection step runs whenever \texttt{l1\_reg} is not \texttt{False}, \texttt{0} or \texttt{'auto'} (\texttt{shap/explainers/\_kernel.py}), so it applies regardless of the sample count. On German Credit ($d{=}20$) the attributions of this arm have ten non-zero entries per instance (checked in the runs). This sparsification makes an error of this size expected, not a discovered failure of SHAP: it is a configuration pitfall of one implementation, and the arm was run at $1\times$ only. The median relative error against the exact target is $\nGerLoneRelErrEstMin$--$\nGerLoneRelErrEstMax$, the top-3 hit rate $\nGerLoneTopHitEstMin$--$\nGerLoneTopHitEstMax$, and all $\nGerLoneVerCond$ cells are classed ``conditional'' (Table~\ref{tab:verdicts}). Repeatability between runs stays high (Spearman $\nGerLoneSpearEstMin$--$\nGerLoneSpearEstMax$) because the bias is shared by all runs, so a repeatability check alone would not detect it; this is why accuracy against an exact target is a separate notion. On ACSEmployment ($d{=}16$) the arm remains within the rule (relative error up to $\nAcsLoneRelErrEstMax$).

\begin{table}[t]
\centering
\caption{Verdicts of the pre-stated decision rule (Sec.~\ref{sec:method}): number of the four models classed reliable, conditional or unreliable in each cell. Cell \texttt{kernel\_q}, \texttt{kernel\_1}, \texttt{kernel\_4}, \texttt{kernel\_16}: Kernel SHAP at $\nicefrac14,1,4,16\times$ the library budget with $\ell_1$ off; \texttt{l1default}: the library default at $1\times$; \texttt{perm}: permutation SHAP. $n{=}50$ instances, $R{=}20$ runs (20 instances at $16\times$). ACS Income and Pima are deterministic-regime controls. Per-model rows and the cells not run are in Appendix~\ref{app:cells}.}
\label{tab:verdicts}
\input{../tables/tab6a_claims_summary.tex}
\end{table}
\noindent The ``four reliable'' counts in Table~\ref{tab:verdicts} include LR (additive, exact) and tree models (exact target); only the MLP is a setting where sampling is the practical option. Under stricter thresholds the number of reliable MLP cells falls from $\nSensMlpBase$ of $\nSensMlpN$ to $\nSensMlpStrictA$ and $\nSensMlpStrictB$ (Appendix~\ref{app:sens}).

\paragraph{Averaging does not beat one larger run.} At equal total evaluations, the mean of $\nEqBudG$ runs at $B/\nEqBudG$ has a mean relative error that is higher than that of one run at $B$ by $\nEqBudRelErrDiffMin$--$\nEqBudRelErrDiffMax$ in absolute terms (non-additive models; four $\nicefrac14\times$ runs against one $1\times$ run; four $1\times$ runs against one $4\times$ run: $\nEqBudFourRelErrDiffMin$--$\nEqBudFourRelErrDiffMax$; bootstrap intervals exclude zero in every comparison); per-model effects and intervals in Table~\ref{tab:eqbud}, while the top-3 hit rates differ by at most $\nEqBudTopHitAbsMax$. The effect is small next to the error levels ($0.004$--$0.02$ at $1\times$), so we read it as ``averaging did not help'', not as a strong negative result. A plausible mechanism, not tested here, is a finite-sample bias of the Kernel SHAP regression that averaging does not remove; permutation SHAP was not compared (no $\nicefrac14\times$ permutation runs).

\begin{figure}[t]
\centering
\includegraphics[width=\linewidth]{fig3_gap_to_noise.pdf}
\caption{Resolvability of the ordering (N3). Fraction of adjacent-rank pairs of the exact target whose attribution gap exceeds $t$ times the run-to-run standard deviation of that gap, as the threshold $t$ varies (dash-dotted: $t=2$). Top: ranks 1 versus 2; bottom: ranks 1--4 pooled. Solid: $1\times$; dotted: $\nicefrac14\times$ budget. Shaded: 95\% instance-bootstrap intervals (German MLP: 20 instances, hence wider). Controls are omitted (every pair is resolvable).}
\label{fig:gap}
\end{figure}

\subsection{Where a single run is less reliable}
\label{sec:conditions}
\paragraph{Instances.} The 50 instances are half random and half near the boundary, so the following are not estimates of a fraction of the test population. Averages hide a tail: on German Credit, $\nGerFracInstBelowNine$ of the instance--model pairs of non-additive models have mean pairwise top-3 Jaccard below $0.9$ at $1\times$ (the three models share instances, so these pairs are not independent), with a minimum of $\nGerMinInstJac$ (Appendix Figure~\ref{fig:instances}); on ACSEmployment the minimum is $\nAcsMinInstJac$. Hit rates against the exact target stay high even for these instances. By stratum, the mean top-3 hit rate is $\nGerHitNear$ (near boundary, 25 instances) against $\nGerHitRand$ (random, 25) on German Credit and $\nAcsHitNear$ against $\nAcsHitRand$ on ACSEmployment; with 25 instances per stratum we do not draw a conclusion about the boundary. Per-model strata and tail shares are in Tables~\ref{tab:strata} and~\ref{tab:tail}, and stricter agreement notions (top-5, ordered top-3) in Table~\ref{tab:stricter}.

\paragraph{Resolvability of the ordering (N3).} The top of the ranking is stable, but fine ordering is not. The fraction of adjacent pairs, in the exact target's ranking, whose gap exceeds twice the run-to-run standard deviation of that gap is $\nGerPairsResMin$--$\nGerPairsResMax$ on German Credit and $\nAcsPairsResMin$--$\nAcsPairsResMax$ on ACSEmployment at $1\times$ (non-additive models); for the first pair (ranks 1 vs 2) it is $\nGerTopPairResMin$--$\nGerTopPairResMax$ and $\nAcsTopPairResMin$--$\nAcsTopPairResMax$, but for pairs at ranks four and below on German Credit only $\nGerLowPairsResMin$--$\nGerLowPairsResMax$. A single run reproduces the exact order of an adjacent pair with probability $\nGerPOrderMin$--$\nGerPOrderMax$ (German Credit) and $\nAcsPOrderMin$--$\nAcsPOrderMax$ (ACSEmployment) on average (Figure~\ref{fig:gap}). The pairs that are not resolved are mostly at lower ranks, where gaps between small attributions are small in any vector, so the order-reproduction probability is the more direct number. Reading a single run's complete ranking therefore over-reports what it supports; the top-3 \emph{set} is supported (Sec.~\ref{sec:main}), the order within it less so.

\paragraph{Noise relative to other variation (N4; descriptive, at a background of 50 rows and $1\times$).} Redrawing the 50-row background changes the attribution vector by $\nGerBgRelSdFacMin$--$\nGerBgRelSdFacMax$\% of its norm on German Credit and $\nAcsBgRelSdFacMin$--$\nAcsBgRelSdFacMax$\% on ACSEmployment (standard deviation over 10 redraws, 3 explainer seeds per redraw), against $\nGerBgRelSdExpMin$--$\nGerBgRelSdExpMax$\% and $\nAcsBgRelSdExpMin$--$\nAcsBgRelSdExpMax$\% from the explainer seed (20 instances; non-additive models). As variances, the background factor is $\nGerBgRatioMin$--$\nGerBgRatioMax$ and $\nAcsBgRatioMin$--$\nAcsBgRatioMax$ times larger, and refitting the model with 10 other seeds gives ratios of $\nGerRefitRatioMin$--$\nGerRefitRatioMax$ and $\nAcsRefitRatioMin$--$\nAcsRefitRatioMax$ (Figure~\ref{fig:var}). A background redraw changes the \emph{target} Shapley vector rather than the estimate of a fixed one, so this ratio shows that the explainer seed is a minor part of the variation a deployed explanation carries; it does not show that the explainer is accurate for some ``correct'' background; refit ratios for LR are undefined (identical models) and the probe has 20 instances. Ratios with intervals per model are in Table~\ref{tab:n4ratio}.

\begin{figure}[t]
\centering
\includegraphics[width=\linewidth]{fig4_variance_components.pdf}
\caption{Noise from three sources (N4): standard deviation of the attribution vector, as a fraction of its norm, across explainer seeds at $1\times$ budget, across 10 redraws of the 50-row background, and across 10 refits of the model (20 instances; the last two with 3 explainer seeds each). Bars: 95\% bootstrap intervals over instances. LR has exact Kernel SHAP, so its explainer noise is drawn at the floor ($10^{-6}$); its refit component is zero by construction and is not drawn. A background redraw changes the target vector itself.}
\label{fig:var}
\end{figure}
